\documentclass[aspectratio=169,11pt]{beamer}
\usetheme{Madrid}
\usecolortheme{whale}
\usepackage[utf8]{inputenc}
\usepackage{graphicx}
\usepackage{amsmath,amssymb}
\usepackage{booktabs}
\usepackage{hyperref}

\title[Talk]{Paper Title}
\subtitle{Subtitle if needed}
\author{Author One \and Author Two}
\institute{University / Institution}
\date{\conference or date}

\begin{document}

\begin{frame}
\titlepage
\end{frame}

\begin{frame}{Outline}
\tableofcontents
\end{frame}

\section{Introduction}

\begin{frame}{Background \& Motivation}
\begin{itemize}
    \item Context: What is the broader area?
    \item Problem: What specific problem does this work address?
    \item Gap: What is missing in existing approaches?
    \item Goal: What does this paper aim to achieve?
\end{itemize}
\end{frame}

\begin{frame}{Contributions}
\begin{enumerate}
    \item \textbf{Contribution 1}: Brief description
    \item \textbf{Contribution 2}: Brief description
    \item \textbf{Contribution 3}: Brief description
\end{enumerate}
\end{frame}

\section{Related Work}

\begin{frame}{Related Work}
\begin{columns}[T]
\begin{column}{0.48\textwidth}
\textbf{Category A:}
\begin{itemize}
    \item Method X (Author, Year): Strength, weakness
    \item Method Y (Author, Year): Strength, weakness
\end{itemize}
\end{column}
\begin{column}{0.48\textwidth}
\textbf{Category B:}
\begin{itemize}
    \item Method Z (Author, Year): Strength, weakness
    \item Method W (Author, Year): Strength, weakness
\end{itemize}
\end{column}
\end{columns}

\vspace{1em}
\textbf{Research Gap:} [State the gap your work addresses]
\end{frame}

\section{Method}

\begin{frame}{Problem Formulation}
\begin{block}{Definition}
    Given $\mathcal{X}$ and $\mathcal{Y}$, learn $f: \mathcal{X} \rightarrow \mathcal{Y}$
\end{block}
\end{frame}

\begin{frame}{Proposed Approach}
\begin{columns}[T]
\begin{column}{0.55\textwidth}
\textbf{Overview:}
\begin{enumerate}
    \item Step 1: Description
    \item Step 2: Description
    \item Step 3: Description
\end{enumerate}

\textbf{Key Innovation:}
Explain the core idea.
\end{column}
\begin{column}{0.42\textwidth}
\includegraphics[width=\textwidth]{method_overview.pdf}
\end{column}
\end{columns}
\end{frame}

\section{Experiments}

\begin{frame}{Experimental Setup}
\begin{columns}[T]
\begin{column}{0.48\textwidth}
\textbf{Datasets:}
\begin{itemize}
    \item Dataset 1: Description (N samples)
    \item Dataset 2: Description (N samples)
\end{itemize}
\end{column}
\begin{column}{0.48\textwidth}
\textbf{Baselines:}
\begin{itemize}
    \item Method A (Year)
    \item Method B (Year)
    \item Method C (Year)
\end{itemize}
\end{column}
\end{columns}
\end{frame}

\begin{frame}{Main Results}
\begin{table}
\centering
\small
\begin{tabular}{lccc}
\toprule
Method & Accuracy & F1 & RMSE \\
\midrule
Baseline A & 0.85 & 0.83 & 0.15 \\
Baseline B & 0.87 & 0.85 & 0.13 \\
Baseline C & 0.89 & 0.87 & 0.11 \\
\textbf{Ours} & \textbf{0.93} & \textbf{0.91} & \textbf{0.08} \\
\bottomrule
\end{tabular}
\caption{Comparison on Dataset 1}
\end{table}
\end{frame}

\begin{frame}{Ablation Study}
\includegraphics[width=0.8\textwidth]{ablation_results.pdf}

\vspace{1em}
\begin{itemize}
    \item Removing component X drops accuracy by Y\%
    \item Component Z is most important for performance
\end{itemize}
\end{frame}

\section{Conclusion}

\begin{frame}{Conclusions \& Future Work}
\textbf{Key Takeaways:}
\begin{enumerate}
    \item Our method achieves state-of-the-art on [task]
    \item The key insight is [core idea]
    \item Analysis shows [finding]
\end{enumerate}

\vspace{1em}
\textbf{Limitations \& Future Work:}
\begin{itemize}
    \item Limitation 1 -> Future direction 1
    \item Limitation 2 -> Future direction 2
\end{itemize}
\end{frame}

\begin{frame}{Thank You}
\centering
\Large Questions?

\vspace{2em}
\small
Contact: email@university.edu

Code: github.com/username/repo
\end{frame}

\end{document}
